\documentclass[a4paper,11pt]{article}

% ---------- Tiếng Việt (pdfLaTeX) ----------
\usepackage[utf8]{inputenc}
\usepackage[T5]{fontenc}
\usepackage[vietnamese]{babel}
\usepackage{lmodern}

% ---------- Các gói theo yêu cầu ----------
\usepackage{amsmath, amssymb, amsfonts}
\usepackage{booktabs}
\usepackage[a4paper,margin=2cm]{geometry}
\usepackage{tikz}
\usepackage{xcolor}
\usepackage{array}
\usetikzlibrary{trees, shapes, arrows, positioning}

% ---------- Các gói bổ trợ ----------
\usepackage{longtable}
\usepackage{graphicx}
\usepackage{enumitem}
\usepackage{float}

\setlength{\parindent}{0pt}
\setlength{\parskip}{4pt}
\renewcommand{\arraystretch}{1.25}

% ---------- Macro ----------
\newcommand{\On}[2]{\mathit{On}(#1,#2)}
\newcommand{\OnT}[1]{\mathit{OnTable}(#1)}
\newcommand{\Clr}[1]{\mathit{Clear}(#1)}
\newcommand{\Hold}[1]{\mathit{Holding}(#1)}
\newcommand{\HE}{\mathit{HandEmpty}}
\newcommand{\Blk}[1]{\mathit{Block}(#1)}
\newcommand{\Op}[2]{\mathsf{#1}(#2)}
\newcommand{\PrA}{\Hold{A}\wedge\Clr{C}}
\newcommand{\PrB}{\OnT{A}\wedge\Clr{A}\wedge\HE}
\newcommand{\PrC}{\On{B}{A}\wedge\Clr{B}\wedge\HE}
\newcommand{\Rst}{\OnT{B}\sm \OnT{C}\sm \Clr{A}\sm \Clr{B}\sm \HE\sm \mathcal{G}}
\newcommand{\cm}{,\allowbreak\ }
\newcommand{\sm}{;\allowbreak\ }
\newcommand{\st}[3]{(#1,#2,#3)}
\newcommand{\ghf}[3]{[#1,#2,#3]}
\newcommand{\nodetxt}[5]{\textbf{#1}\ $(#2)$\\ $g{=}#3\cm h{=}#4\cm f{=}#5$}
\newcommand{\cau}[1]{\section*{Câu #1}\addcontentsline{toc}{section}{Câu #1}}

\title{\textbf{BÀI ÔN TẬP TRÍ TUỆ NHÂN TẠO NÂNG CAO}\\[2pt]\large Lời giải chi tiết}
\author{}
\date{}

\begin{document}
\maketitle

%====================================================================
\cau{1}
\textbf{Đề:} Ba bình A, B, C có dung tích $8$, $5$, $3$ lít; ban đầu A đầy, B và C rỗng; chỉ được rót qua lại giữa các bình. Cần đạt trạng thái có đúng $4$ lít ở hai bình bất kỳ.

\subsection*{a) Mô hình không gian trạng thái}

\paragraph{Biểu diễn trạng thái.} Mỗi trạng thái là một bộ ba
\[
s=(x_A,x_B,x_C),\qquad 0\le x_A\le 8\cm 0\le x_B\le 5\cm 0\le x_C\le 3,
\]
trong đó $x_X$ là lượng nước (lít) trong bình $X$. Vì không có nguồn nước bổ sung và không có nước bị đổ đi nên luôn có bất biến
$x_A+x_B+x_C=8$. Gọi $c=(c_A,c_B,c_C)=(8,5,3)$ là vector dung tích.

\paragraph{Trạng thái bắt đầu.} $s_0=(8,0,0)$.

\paragraph{Trạng thái kết thúc.} Cần có hai bình chứa đúng $4$ lít. Do tổng luôn bằng $8$ nên bình còn lại phải rỗng ($4+4+0=8$). Bình C chỉ chứa tối đa $3<4$ lít nên không thể là một trong hai bình chứa $4$ lít; do đó hai bình đó chỉ có thể là A và B:
\[
\mathrm{Goal}=\{(4,4,0)\}.
\]
(Các bộ $(4,0,4)$ và $(0,4,4)$ không hợp lệ vì vi phạm $x_C\le 3$.)

\paragraph{Các toán tử (luật chuyển trạng thái).} Có $6$ toán tử rót $X\to Y$ ($X\neq Y$, $X,Y\in\{A,B,C\}$): rót nước từ bình $X$ sang bình $Y$ cho đến khi $X$ cạn hoặc $Y$ đầy.
\begin{itemize}[nosep]
  \item \textbf{Điều kiện áp dụng:} $x_X>0$ và $x_Y<c_Y$.
  \item \textbf{Lượng nước rót:} $q=\min\bigl(x_X\cm c_Y-x_Y\bigr)$.
  \item \textbf{Trạng thái mới:} $x_X'=x_X-q$, $x_Y'=x_Y+q$, bình còn lại giữ nguyên.
  \item \textbf{Chi phí} mỗi bước rót: $1$.
\end{itemize}

\begin{center}
\begin{tabular}{@{}lll@{}}
\toprule
Toán tử & Điều kiện & Tác động \\
\midrule
$A\to B$ & $x_A>0\cm x_B<5$ & $q=\min(x_A,5-x_B)$\sm $x_A\mathrel{-}=q\cm x_B\mathrel{+}=q$ \\
$A\to C$ & $x_A>0\cm x_C<3$ & $q=\min(x_A,3-x_C)$\sm $x_A\mathrel{-}=q\cm x_C\mathrel{+}=q$ \\
$B\to A$ & $x_B>0\cm x_A<8$ & $q=\min(x_B,8-x_A)$\sm $x_B\mathrel{-}=q\cm x_A\mathrel{+}=q$ \\
$B\to C$ & $x_B>0\cm x_C<3$ & $q=\min(x_B,3-x_C)$\sm $x_B\mathrel{-}=q\cm x_C\mathrel{+}=q$ \\
$C\to A$ & $x_C>0\cm x_A<8$ & $q=\min(x_C,8-x_A)$\sm $x_C\mathrel{-}=q\cm x_A\mathrel{+}=q$ \\
$C\to B$ & $x_C>0\cm x_B<5$ & $q=\min(x_C,5-x_B)$\sm $x_C\mathrel{-}=q\cm x_B\mathrel{+}=q$ \\
\bottomrule
\end{tabular}
\end{center}

\subsection*{b) Thuật toán A* với $h(x_A,x_B,x_C)=|x_A-4|+|x_B-4|+|x_C-4|$}

\paragraph{Quy ước.} $g(n)$ = số lần rót từ $s_0$ đến $n$; $f(n)=g(n)+h(n)$. Mỗi bước chọn nút trong OPEN có $f$ nhỏ nhất; nếu bằng nhau thì chọn nút có $h$ nhỏ hơn, nếu vẫn bằng nhau thì chọn nút sinh ra trước (thứ tự toán tử: $A\!\to\! B\cm A\!\to\! C\cm B\!\to\! A\cm B\!\to\! C\cm C\!\to\! A\cm C\!\to\! B$). Nút đã nằm trong CLOSED thì bỏ; nút đã nằm trong OPEN với $g$ nhỏ hơn hoặc bằng thì bỏ bản sao mới. Kiểm tra đích khi \emph{mở rộng} nút. Ký hiệu $\ghf{g}{h}{f}$.

\paragraph{Nhận xét.} $h(8,0,0)=4+4+4=12$. Với trạng thái đích $h(4,4,0)=0+0+4=4\neq 0$, và $h$ không chấp nhận được (ví dụ $h(1,4,3)=4$ nhưng chỉ cần $1$ bước rót để tới đích), nên A* ở đây chỉ được dùng như tìm kiếm heuristic; kết quả vẫn tìm được lời giải $7$ bước.

\paragraph{Các bước duyệt nút.}
\footnotesize
\begin{longtable}{@{}>{\centering\arraybackslash}p{0.55cm} >{\raggedright\arraybackslash}p{2.6cm} >{\raggedright\arraybackslash}p{6.9cm} >{\raggedright\arraybackslash}p{5.0cm}@{}}
\toprule
\textbf{Bước} & \textbf{Nút mở rộng} $\ghf{g}{h}{f}$ & \textbf{Các nút con} (toán tử: trạng thái $\ghf{g}{h}{f}$) & \textbf{OPEN sau bước} (trạng thái:$f$) \\
\midrule
\endhead
1 & $\st800\ \ghf{0}{12}{12}$ & $A\to B$: $\st350\ \ghf{1}{6}{7}$ \newline $A\to C$: $\st503\ \ghf{1}{6}{7}$ & $\st350{:}7\cm \st503{:}7$ \\
2 & $\st350\ \ghf{1}{6}{7}$ & $A\to C$: $\st053\ \ghf{2}{6}{8}$ \newline $B\to A$: $\st800$ (CLOSED, bỏ) \newline $B\to C$: $\st323\ \ghf{2}{4}{6}$ & $\st323{:}6\cm \st503{:}7\cm \st053{:}8$ \\
3 & $\st323\ \ghf{2}{4}{6}$ & $A\to B$: $\st053\ \ghf{3}{6}{9}$ (đã có $g{=}2$, bỏ) \newline $B\to A$: $\st503\ \ghf{3}{6}{9}$ (đã có $g{=}1$, bỏ) \newline $C\to A$: $\st620\ \ghf{3}{8}{11}$ \newline $C\to B$: $\st350$ (CLOSED, bỏ) & $\st503{:}7\cm \st053{:}8\cm \st620{:}11$ \\
4 & $\st503\ \ghf{1}{6}{7}$ & $A\to B$: $\st053\ \ghf{2}{6}{8}$ (trùng, bỏ) \newline $C\to A$: $\st800$ (CLOSED, bỏ) \newline $C\to B$: $\st530\ \ghf{2}{6}{8}$ & $\st053{:}8\cm \st530{:}8\cm \st620{:}11$ \\
5 & $\st053\ \ghf{2}{6}{8}$ & $B\to A$: $\st503$ (CLOSED) \newline $C\to A$: $\st350$ (CLOSED) $\Rightarrow$ \emph{nhánh cụt} & $\st530{:}8\cm \st620{:}11$ \\
6 & $\st530\ \ghf{2}{6}{8}$ & $A\to B$: $\st350$ (CLOSED) \newline $A\to C$: $\st233\ \ghf{3}{4}{7}$ \newline $B\to A$: $\st800$ (CLOSED) \newline $B\to C$: $\st503$ (CLOSED) & $\st233{:}7\cm \st620{:}11$ \\
7 & $\st233\ \ghf{3}{4}{7}$ & $A\to B$: $\st053$, $B\to A$: $\st503$, $C\to A$: $\st530$ (đều CLOSED) \newline $C\to B$: $\st251\ \ghf{4}{6}{10}$ & $\st251{:}10\cm \st620{:}11$ \\
8 & $\st251\ \ghf{4}{6}{10}$ & $A\to C$: $\st053$ (CLOSED) \newline $B\to A$: $\st701\ \ghf{5}{10}{15}$ \newline $B\to C$: $\st233$, $C\to A$: $\st350$ (CLOSED) & $\st620{:}11\cm \st701{:}15$ \\
9 & $\st620\ \ghf{3}{8}{11}$ & $A\to B$: $\st350$, $A\to C$: $\st323$, $B\to A$: $\st800$ (CLOSED) \newline $B\to C$: $\st602\ \ghf{4}{8}{12}$ & $\st602{:}12\cm \st701{:}15$ \\
10 & $\st602\ \ghf{4}{8}{12}$ & $A\to B$: $\st152\ \ghf{5}{6}{11}$ \newline $A\to C$: $\st503$, $C\to A$: $\st800$, $C\to B$: $\st620$ (CLOSED) & $\st152{:}11\cm \st701{:}15$ \\
11 & $\st152\ \ghf{5}{6}{11}$ & $B\to C$: $\st143\ \ghf{6}{4}{10}$ \newline $A\to C$: $\st053$, $B\to A$: $\st602$, $C\to A$: $\st350$ (CLOSED) & $\st143{:}10\cm \st701{:}15$ \\
12 & $\st143\ \ghf{6}{4}{10}$ & $C\to A$: $\st440\ \ghf{7}{4}{11}$ \newline $A\to B$: $\st053$, $B\to A$: $\st503$, $C\to B$: $\st152$ (CLOSED) & $\st440{:}11\cm \st701{:}15$ \\
13 & $\st440\ \ghf{7}{4}{11}$ & \textbf{Là trạng thái đích $\Rightarrow$ DỪNG.} & --- \\
\bottomrule
\end{longtable}
\normalsize

\paragraph{Lời giải tìm được} ($7$ bước, chi phí $g=7$):
\begin{gather*}
\st800\xrightarrow{A\to B}\st350\xrightarrow{B\to C}\st323\xrightarrow{C\to A}\st620\xrightarrow{B\to C}\st602\\
\xrightarrow{A\to B}\st152\xrightarrow{B\to C}\st143\xrightarrow{C\to A}\st440.
\end{gather*}

\paragraph{Cây tìm kiếm A*.} Số trong ngoặc \textbf{[k]} là thứ tự mở rộng nút; nút tô xanh nằm trên đường đi lời giải, nút tô đỏ là đích, nét đứt là nút còn trong OPEN chưa được mở rộng, các nút lặp (CLOSED) không vẽ lại.

\begin{figure}[H]
\centering
\begin{tikzpicture}[x=1cm,y=1cm,
  nd/.style={draw,rounded corners=2pt,align=center,font=\scriptsize,inner sep=2.5pt,minimum width=2.3cm},
  sol/.style={nd,fill=green!15,thick},
  goalnd/.style={nd,fill=red!20,very thick},
  openn/.style={nd,dashed,fill=gray!10},
  dead/.style={nd,fill=yellow!20},
  lab/.style={midway,fill=white,inner sep=1pt,font=\tiny}]
\node[sol]    (n1)  at (0,0)       {\nodetxt{[1]}{8,0,0}{0}{12}{12}};
\node[sol]    (n2)  at (-4.9,-1.9) {\nodetxt{[2]}{3,5,0}{1}{6}{7}};
\node[nd]     (n4)  at (4.2,-1.9)  {\nodetxt{[4]}{5,0,3}{1}{6}{7}};
\node[dead]   (n5)  at (-7.0,-3.8) {\nodetxt{[5]}{0,5,3}{2}{6}{8}};
\node[sol]    (n3)  at (-3.0,-3.8) {\nodetxt{[3]}{3,2,3}{2}{4}{6}};
\node[sol]    (n9)  at (-3.0,-5.7) {\nodetxt{[9]}{6,2,0}{3}{8}{11}};
\node[sol]    (n10) at (-3.0,-7.6) {\nodetxt{[10]}{6,0,2}{4}{8}{12}};
\node[sol]    (n11) at (-3.0,-9.5) {\nodetxt{[11]}{1,5,2}{5}{6}{11}};
\node[sol]    (n12) at (-3.0,-11.4){\nodetxt{[12]}{1,4,3}{6}{4}{10}};
\node[goalnd] (n13) at (-3.0,-13.3){\nodetxt{[13]}{4,4,0}{7}{4}{11}};
\node[nd]     (n6)  at (4.2,-3.8)  {\nodetxt{[6]}{5,3,0}{2}{6}{8}};
\node[nd]     (n7)  at (4.2,-5.7)  {\nodetxt{[7]}{2,3,3}{3}{4}{7}};
\node[nd]     (n8)  at (4.2,-7.6)  {\nodetxt{[8]}{2,5,1}{4}{6}{10}};
\node[openn]  (n14) at (4.2,-9.5)  {\nodetxt{OPEN}{7,0,1}{5}{10}{15}};
\draw[->,very thick,green!50!black] (n1)  -- node[lab]{$A\!\to\! B$} (n2);
\draw[->] (n1)  -- node[lab]{$A\!\to\! C$} (n4);
\draw[->] (n2)  -- node[lab]{$A\!\to\! C$} (n5);
\draw[->,very thick,green!50!black] (n2)  -- node[lab]{$B\!\to\! C$} (n3);
\draw[->,very thick,green!50!black] (n3)  -- node[lab]{$C\!\to\! A$} (n9);
\draw[->,very thick,green!50!black] (n9)  -- node[lab]{$B\!\to\! C$} (n10);
\draw[->,very thick,green!50!black] (n10) -- node[lab]{$A\!\to\! B$} (n11);
\draw[->,very thick,green!50!black] (n11) -- node[lab]{$B\!\to\! C$} (n12);
\draw[->,very thick,green!50!black] (n12) -- node[lab]{$C\!\to\! A$} (n13);
\draw[->] (n4)  -- node[lab]{$C\!\to\! B$} (n6);
\draw[->] (n6)  -- node[lab]{$A\!\to\! C$} (n7);
\draw[->] (n7)  -- node[lab]{$C\!\to\! B$} (n8);
\draw[->,dashed] (n8)  -- node[lab]{$B\!\to\! A$} (n14);
\end{tikzpicture}
\caption{Cây tìm kiếm A* của bài toán đong nước (đích $(4,4,0)$, $g=7$).}
\end{figure}

%====================================================================
\clearpage
\cau{2}
\textbf{Đề:} Cây trò chơi $2$ người (gốc MAX $\to$ MIN $\to$ MAX $\to$ MIN); người chơi chọn nước đi có giá trị lượng giá càng lớn càng tốt. Để tiện tham chiếu, đặt tên các nút trong:
gốc $A$ (MAX); $B_1,B_2,B_3$ (MIN); $C_1,\dots,C_5$ (MAX); $D_1,\dots,D_4$ (MIN). Cấu trúc cây:
\begin{itemize}[nosep]
  \item $B_1$: các con $C_1\cm C_2\cm \text{lá }11$; \quad $C_1$: lá $4,13$; \quad $C_2$: $D_1$ và lá $15$; \quad $D_1$: lá $5,10$.
  \item $B_2$: các con $C_3\cm \text{lá }12$; \quad $C_3$: $D_2$, lá $1$, $D_3$; \quad $D_2$: lá $9,8$; \quad $D_3$: lá $6,12$.
  \item $B_3$: các con $C_4\cm C_5$; \quad $C_4$: lá $10,3$ và $D_4$; \quad $D_4$: lá $7,5,2$; \quad $C_5$: lá $7,9$.
\end{itemize}

\subsection*{a) Giải thuật Minimax}
Tính từ dưới lên: nút MAX lấy giá trị lớn nhất của các con, nút MIN lấy giá trị nhỏ nhất.
\begin{align*}
D_1&=\min(5,10)=5, & D_2&=\min(9,8)=8, & D_3&=\min(6,12)=6, \\
D_4&=\min(7,5,2)=2, & C_1&=\max(4,13)=13, & C_2&=\max(5,15)=15, \\
C_3&=\max(8,1,6)=8, & C_4&=\max(10,3,2)=10, & C_5&=\max(7,9)=9,\\
B_1&=\min(13,15,11)=11, & B_2&=\min(8,12)=8, & B_3&=\min(10,9)=9,\\
A&=\max(11,8,9)=11.
\end{align*}
Vậy giá trị của gốc là $11$; nước đi tối ưu của MAX là đi tới $B_1$ (rồi MIN đi tới lá $11$).

\begin{figure}[H]
\centering
\resizebox{\textwidth}{!}{%
\begin{tikzpicture}[x=1cm,y=1cm,
  maxn/.style={circle,draw,thick,fill=green!15,minimum size=0.72cm,inner sep=0pt,font=\small},
  minn/.style={rectangle,draw,thick,fill=yellow!25,minimum size=0.66cm,inner sep=0pt,font=\small},
  leafn/.style={circle,draw,fill=white,minimum size=0.66cm,inner sep=0pt,font=\small},
  cut/.style={circle,draw=gray,dashed,fill=gray!15,text=gray,minimum size=0.66cm,inner sep=0pt,font=\small},
  ablab/.style={font=\tiny,inner sep=1pt,fill=white,fill opacity=0.85,text opacity=1}]
\node[maxn] (n0) at (9.58,0.00) {11};
\node[minn] (n1) at (2.89,-1.75) {11};
\node[maxn] (n2) at (0.53,-3.50) {13};
\node[leafn] (n3) at (0.00,-5.25) {4};
\node[leafn] (n4) at (1.05,-5.25) {13};
\node[maxn] (n5) at (3.41,-3.50) {15};
\node[minn] (n6) at (2.62,-5.25) {5};
\node[leafn] (n7) at (2.10,-7.00) {5};
\node[leafn] (n8) at (3.15,-7.00) {10};
\node[leafn] (n9) at (4.20,-5.25) {15};
\node[leafn] (n10) at (5.25,-3.50) {11};
\node[minn] (n11) at (9.98,-1.75) {8};
\node[maxn] (n12) at (8.40,-3.50) {8};
\node[minn] (n13) at (6.82,-5.25) {8};
\node[leafn] (n14) at (6.30,-7.00) {9};
\node[leafn] (n15) at (7.35,-7.00) {8};
\node[leafn] (n16) at (8.40,-5.25) {1};
\node[minn] (n17) at (9.98,-5.25) {6};
\node[leafn] (n18) at (9.45,-7.00) {6};
\node[leafn] (n19) at (10.50,-7.00) {12};
\node[leafn] (n20) at (11.55,-3.50) {12};
\node[minn] (n21) at (16.28,-1.75) {9};
\node[maxn] (n22) at (14.18,-3.50) {10};
\node[leafn] (n23) at (12.60,-5.25) {10};
\node[leafn] (n24) at (13.65,-5.25) {3};
\node[minn] (n25) at (15.75,-5.25) {2};
\node[leafn] (n26) at (14.70,-7.00) {7};
\node[leafn] (n27) at (15.75,-7.00) {5};
\node[leafn] (n28) at (16.80,-7.00) {2};
\node[maxn] (n29) at (18.38,-3.50) {9};
\node[leafn] (n30) at (17.85,-5.25) {7};
\node[leafn] (n31) at (18.90,-5.25) {9};
\draw[red,very thick] (n0) -- (n1);
\draw (n0) -- (n11);
\draw (n0) -- (n21);
\draw (n1) -- (n2);
\draw (n1) -- (n5);
\draw[red,very thick] (n1) -- (n10);
\draw (n2) -- (n3);
\draw (n2) -- (n4);
\draw (n5) -- (n6);
\draw (n5) -- (n9);
\draw (n6) -- (n7);
\draw (n6) -- (n8);
\draw (n11) -- (n12);
\draw (n11) -- (n20);
\draw (n12) -- (n13);
\draw (n12) -- (n16);
\draw (n12) -- (n17);
\draw (n13) -- (n14);
\draw (n13) -- (n15);
\draw (n17) -- (n18);
\draw (n17) -- (n19);
\draw (n21) -- (n22);
\draw (n21) -- (n29);
\draw (n22) -- (n23);
\draw (n22) -- (n24);
\draw (n22) -- (n25);
\draw (n25) -- (n26);
\draw (n25) -- (n27);
\draw (n25) -- (n28);
\draw (n29) -- (n30);
\draw (n29) -- (n31);
\node[font=\tiny\bfseries,text=blue!60!black,left=-1pt of n0] {$A$};
\node[font=\tiny\bfseries,text=blue!60!black,left=-1pt of n1] {$B_1$};
\node[font=\tiny\bfseries,text=blue!60!black,left=-1pt of n2] {$C_1$};
\node[font=\tiny\bfseries,text=blue!60!black,left=-1pt of n5] {$C_2$};
\node[font=\tiny\bfseries,text=blue!60!black,left=-1pt of n6] {$D_1$};
\node[font=\tiny\bfseries,text=blue!60!black,left=-1pt of n11] {$B_2$};
\node[font=\tiny\bfseries,text=blue!60!black,left=-1pt of n12] {$C_3$};
\node[font=\tiny\bfseries,text=blue!60!black,left=-1pt of n13] {$D_2$};
\node[font=\tiny\bfseries,text=blue!60!black,left=-1pt of n17] {$D_3$};
\node[font=\tiny\bfseries,text=blue!60!black,left=-1pt of n21] {$B_3$};
\node[font=\tiny\bfseries,text=blue!60!black,left=-1pt of n22] {$C_4$};
\node[font=\tiny\bfseries,text=blue!60!black,left=-1pt of n25] {$D_4$};
\node[font=\tiny\bfseries,text=blue!60!black,left=-1pt of n29] {$C_5$};
\node[font=\bfseries\small,anchor=east] at (-1.60,0.00) {MAX};
\node[font=\bfseries\small,anchor=east] at (-1.60,-1.75) {MIN};
\node[font=\bfseries\small,anchor=east] at (-1.60,-3.50) {MAX};
\node[font=\bfseries\small,anchor=east] at (-1.60,-5.25) {MIN};
\end{tikzpicture}
}
\caption{Cây Minimax với đầy đủ giá trị tại mọi nút (đường đỏ: đường đi tối ưu; hình tròn = MAX, hình vuông = MIN).}
\end{figure}

\subsection*{b) Cắt tỉa Alpha-Beta (duyệt từ trái sang phải)}
Khởi tạo $[\alpha,\beta]=[-\infty,+\infty]$ tại gốc. Nút MAX cập nhật $\alpha=\max(\alpha,v)$, nút MIN cập nhật $\beta=\min(\beta,v)$; \textbf{cắt tỉa} các con còn lại của một nút ngay khi $\alpha\ge\beta$. Khoảng $[\alpha,\beta]$ ghi dưới đây là khoảng \emph{khi vào nút} (truyền từ cha xuống).

\footnotesize
\begin{longtable}{@{}p{1.0cm} >{\raggedright\arraybackslash}p{2.5cm} >{\raggedright\arraybackslash}p{8.3cm} >{\raggedright\arraybackslash}p{3.3cm}@{}}
\toprule
\textbf{Nút} & $[\alpha,\beta]$ \textbf{khi vào} & \textbf{Diễn biến} & \textbf{Kết quả} \\
\midrule
\endhead
$A$ & $[-\infty,+\infty]$ & Gọi lần lượt $B_1,B_2,B_3$. & --- \\
$B_1$ & $[-\infty,+\infty]$ & Thăm $C_1$, $C_2$, lá $11$. & --- \\
$C_1$ & $[-\infty,+\infty]$ & Lá $4$: $\alpha=4$; lá $13$: $\alpha=13$. & trả về $13$ \\
 & & Về $B_1$: $\beta=\min(+\infty,13)=13$. & \\
$C_2$ & $[-\infty,13]$ & Vào $D_1$. & --- \\
$D_1$ & $[-\infty,13]$ & Lá $5$: $\beta=5$ (vì $\alpha=-\infty<5$ nên không cắt); lá $10$: $\beta$ vẫn là $5$. & trả về $5$ \\
 & & Về $C_2$: $\alpha=5$; lá $15$: $v=15$, $\alpha=15\ge\beta=13$ nhưng $15$ là con cuối $\Rightarrow$ không có gì để cắt. & $C_2$ trả về $15$ \\
 & & Về $B_1$: $\beta=\min(13,15)=13$; lá $11$: $\beta=11$. Hết con. & $B_1$ trả về $11$ \\
$A$ & & $\alpha=\max(-\infty,11)=11$. & $[\alpha,\beta]=[11,+\infty]$ \\
\midrule
$B_2$ & $[11,+\infty]$ & Thăm $C_3$ trước. & --- \\
$C_3$ & $[11,+\infty]$ & Vào $D_2$. & --- \\
$D_2$ & $[11,+\infty]$ & Lá $9$: $\beta=9\le\alpha=11$ $\Rightarrow$ \textbf{CẮT} lá $8$. & trả về $9$ \\
 & & Về $C_3$: $v=9$, $\alpha=\max(11,9)=11$; lá $1$: $v$ vẫn là $9$. & \\
$D_3$ & $[11,+\infty]$ & Lá $6$: $\beta=6\le\alpha=11$ $\Rightarrow$ \textbf{CẮT} lá $12$ (con của $D_3$). & trả về $6$ \\
 & & Về $C_3$: $v=\max(9,1,6)=9$. Hết con. & $C_3$ trả về $9$ \\
 & & Về $B_2$: $\beta=\min(+\infty,9)=9\le\alpha=11$ $\Rightarrow$ \textbf{CẮT} lá $12$ (con phải của $B_2$). & $B_2$ trả về $9$ ($\le 11$) \\
 & & (Giá trị thật của $B_2$ là $8$ nhưng $B_2\le 9<11$ nên không thể được MAX chọn.) & \\
$A$ & & $\alpha=\max(11,9)=11$ (không đổi). & \\
\midrule
$B_3$ & $[11,+\infty]$ & Thăm $C_4$ trước. & --- \\
$C_4$ & $[11,+\infty]$ & Lá $10$: $v=10$; lá $3$: $v=10$; vào $D_4$. & --- \\
$D_4$ & $[11,+\infty]$ & Lá $7$: $\beta=7\le\alpha=11$ $\Rightarrow$ \textbf{CẮT} lá $5$ và lá $2$. & trả về $7$ \\
 & & Về $C_4$: $v=\max(10,3,7)=10$. & $C_4$ trả về $10$ \\
 & & Về $B_3$: $\beta=\min(+\infty,10)=10\le\alpha=11$ $\Rightarrow$ \textbf{CẮT} toàn bộ nhánh $C_5$ (gồm lá $7$ và $9$). & $B_3$ trả về $10$ ($\le 11$) \\
$A$ & & $\alpha=\max(11,10)=11$. Hết con. & $A=11$ \\
\bottomrule
\end{longtable}
\normalsize

\paragraph{Kết luận.} Giá trị gốc $=11$ (giống Minimax), nước đi tốt nhất là tới $B_1$. Các nhánh bị cắt tỉa:
\begin{enumerate}[nosep]
  \item lá $8$ (con phải của $D_2$), do $\beta=9\le\alpha=11$;
  \item lá $12$ (con phải của $D_3$), do $\beta=6\le\alpha=11$;
  \item lá $12$ (con phải của $B_2$), do $\beta=9\le\alpha=11$;
  \item lá $5$ và lá $2$ (con của $D_4$), do $\beta=7\le\alpha=11$;
  \item toàn bộ nhánh $C_5$ cùng hai lá $7,9$ (con phải của $B_3$), do $\beta=10\le\alpha=11$.
\end{enumerate}
Tổng cộng cắt $7$ trong $19$ lá; chỉ đánh giá $12$ lá ($4,13,5,10,15,11,9,1,6,10,3,7$).

\begin{figure}[H]
\centering
\resizebox{\textwidth}{!}{%
\begin{tikzpicture}[x=1cm,y=1cm,
  maxn/.style={circle,draw,thick,fill=green!15,minimum size=0.72cm,inner sep=0pt,font=\small},
  minn/.style={rectangle,draw,thick,fill=yellow!25,minimum size=0.66cm,inner sep=0pt,font=\small},
  leafn/.style={circle,draw,fill=white,minimum size=0.66cm,inner sep=0pt,font=\small},
  cut/.style={circle,draw=gray,dashed,fill=gray!15,text=gray,minimum size=0.66cm,inner sep=0pt,font=\small},
  ablab/.style={font=\tiny,inner sep=1pt,fill=white,fill opacity=0.85,text opacity=1}]
\node[maxn] (n0) at (9.58,0.00) {11};
\node[minn] (n1) at (2.89,-1.75) {11};
\node[maxn] (n2) at (0.53,-3.50) {13};
\node[leafn] (n3) at (0.00,-5.25) {4};
\node[leafn] (n4) at (1.05,-5.25) {13};
\node[maxn] (n5) at (3.41,-3.50) {15};
\node[minn] (n6) at (2.62,-5.25) {5};
\node[leafn] (n7) at (2.10,-7.00) {5};
\node[leafn] (n8) at (3.15,-7.00) {10};
\node[leafn] (n9) at (4.20,-5.25) {15};
\node[leafn] (n10) at (5.25,-3.50) {11};
\node[minn] (n11) at (9.98,-1.75) {9};
\node[maxn] (n12) at (8.40,-3.50) {9};
\node[minn] (n13) at (6.82,-5.25) {9};
\node[leafn] (n14) at (6.30,-7.00) {9};
\node[cut] (n15) at (7.35,-7.00) {8};
\node[leafn] (n16) at (8.40,-5.25) {1};
\node[minn] (n17) at (9.98,-5.25) {6};
\node[leafn] (n18) at (9.45,-7.00) {6};
\node[cut] (n19) at (10.50,-7.00) {12};
\node[cut] (n20) at (11.55,-3.50) {12};
\node[minn] (n21) at (16.28,-1.75) {10};
\node[maxn] (n22) at (14.18,-3.50) {10};
\node[leafn] (n23) at (12.60,-5.25) {10};
\node[leafn] (n24) at (13.65,-5.25) {3};
\node[minn] (n25) at (15.75,-5.25) {7};
\node[leafn] (n26) at (14.70,-7.00) {7};
\node[cut] (n27) at (15.75,-7.00) {5};
\node[cut] (n28) at (16.80,-7.00) {2};
\node[cut] (n29) at (18.38,-3.50) {?};
\node[cut] (n30) at (17.85,-5.25) {7};
\node[cut] (n31) at (18.90,-5.25) {9};
\draw (n0) -- (n1);
\draw (n0) -- (n11);
\draw (n0) -- (n21);
\draw (n1) -- (n2);
\draw (n1) -- (n5);
\draw (n1) -- (n10);
\draw (n2) -- (n3);
\draw (n2) -- (n4);
\draw (n5) -- (n6);
\draw (n5) -- (n9);
\draw (n6) -- (n7);
\draw (n6) -- (n8);
\draw (n11) -- (n12);
\draw[gray,dashed] (n11) -- (n20);
\draw[red,very thick] (10.63,-2.75)--(10.89,-2.50) (10.63,-2.50)--(10.89,-2.75);
\draw (n12) -- (n13);
\draw (n12) -- (n16);
\draw (n12) -- (n17);
\draw (n13) -- (n14);
\draw[gray,dashed] (n13) -- (n15);
\draw[red,very thick] (6.96,-6.25)--(7.22,-6.00) (6.96,-6.00)--(7.22,-6.25);
\draw (n17) -- (n18);
\draw[gray,dashed] (n17) -- (n19);
\draw[red,very thick] (10.11,-6.25)--(10.37,-6.00) (10.11,-6.00)--(10.37,-6.25);
\draw (n21) -- (n22);
\draw[gray,dashed] (n21) -- (n29);
\draw[red,very thick] (17.20,-2.75)--(17.46,-2.50) (17.20,-2.50)--(17.46,-2.75);
\draw (n22) -- (n23);
\draw (n22) -- (n24);
\draw (n22) -- (n25);
\draw (n25) -- (n26);
\draw[gray,dashed] (n25) -- (n27);
\draw[red,very thick] (15.62,-6.25)--(15.88,-6.00) (15.62,-6.00)--(15.88,-6.25);
\draw[gray,dashed] (n25) -- (n28);
\draw[red,very thick] (16.15,-6.25)--(16.41,-6.00) (16.15,-6.00)--(16.41,-6.25);
\draw[gray,dashed] (n29) -- (n30);
\draw[gray,dashed] (n29) -- (n31);
\node[ablab,above right=-2pt and -3pt of n0] {$[-\infty,+\infty]$};
\node[ablab,above right=-2pt and -3pt of n1] {$[-\infty,+\infty]$};
\node[ablab,above right=-2pt and -3pt of n2] {$[-\infty,+\infty]$};
\node[ablab,above right=-2pt and -3pt of n5] {$[-\infty,13]$};
\node[ablab,above right=-2pt and -3pt of n6] {$[-\infty,13]$};
\node[ablab,above right=-2pt and -3pt of n11] {$[11,+\infty]$};
\node[ablab,above right=-2pt and -3pt of n12] {$[11,+\infty]$};
\node[ablab,above right=-2pt and -3pt of n13] {$[11,+\infty]$};
\node[ablab,above right=-2pt and -3pt of n17] {$[11,+\infty]$};
\node[ablab,above right=-2pt and -3pt of n21] {$[11,+\infty]$};
\node[ablab,above right=-2pt and -3pt of n22] {$[11,+\infty]$};
\node[ablab,above right=-2pt and -3pt of n25] {$[11,+\infty]$};
\node[font=\tiny\bfseries,text=blue!60!black,left=-1pt of n0] {$A$};
\node[font=\tiny\bfseries,text=blue!60!black,left=-1pt of n1] {$B_1$};
\node[font=\tiny\bfseries,text=blue!60!black,left=-1pt of n2] {$C_1$};
\node[font=\tiny\bfseries,text=blue!60!black,left=-1pt of n5] {$C_2$};
\node[font=\tiny\bfseries,text=blue!60!black,left=-1pt of n6] {$D_1$};
\node[font=\tiny\bfseries,text=blue!60!black,left=-1pt of n11] {$B_2$};
\node[font=\tiny\bfseries,text=blue!60!black,left=-1pt of n12] {$C_3$};
\node[font=\tiny\bfseries,text=blue!60!black,left=-1pt of n13] {$D_2$};
\node[font=\tiny\bfseries,text=blue!60!black,left=-1pt of n17] {$D_3$};
\node[font=\tiny\bfseries,text=blue!60!black,left=-1pt of n21] {$B_3$};
\node[font=\tiny\bfseries,text=blue!60!black,left=-1pt of n22] {$C_4$};
\node[font=\tiny\bfseries,text=blue!60!black,left=-1pt of n25] {$D_4$};
\node[font=\tiny\bfseries,text=blue!60!black,left=-1pt of n29] {$C_5$};
\node[font=\bfseries\small,anchor=east] at (-1.60,0.00) {MAX};
\node[font=\bfseries\small,anchor=east] at (-1.60,-1.75) {MIN};
\node[font=\bfseries\small,anchor=east] at (-1.60,-3.50) {MAX};
\node[font=\bfseries\small,anchor=east] at (-1.60,-5.25) {MIN};
\end{tikzpicture}
}
\caption{Cây Alpha-Beta: giá trị cạnh nút là giá trị trả về, $[\alpha,\beta]$ là khoảng khi vào nút; dấu $\times$ đỏ đánh dấu nhánh bị cắt, nút xám nét đứt là phần không được thăm.}
\end{figure}

%====================================================================
\clearpage
\cau{3}
\textbf{Đề:} Mạng Bayes Cloudy ($C$) $\to$ Sprinkler ($S$), Rain ($R$) $\to$ WetGrass ($W$).

\begin{figure}[H]
\centering
\resizebox{0.97\textwidth}{!}{%
\begin{tikzpicture}[
  bn/.style={ellipse,draw,thick,fill=blue!8,minimum width=2.3cm,minimum height=0.9cm,font=\bfseries\small},
  cpt/.style={draw,fill=white,font=\scriptsize,inner sep=2pt,rounded corners=2pt},
  >=latex]
\node[bn] (C) at (0,0) {Cloudy};
\node[bn] (S) at (-3.6,-2.4) {Sprinkler};
\node[bn] (R) at (3.6,-2.4) {Rain};
\node[bn] (W) at (0,-4.8) {WetGrass};
\draw[->,thick] (C) -- (S);
\draw[->,thick] (C) -- (R);
\draw[->,thick] (S) -- (W);
\draw[->,thick] (R) -- (W);
\node[cpt,anchor=east] at (-1.6,0.1) {\begin{tabular}{cc}$P(C{=}T)$&$P(C{=}F)$\\\hline 0.5&0.5\end{tabular}};
\node[cpt,anchor=south west] at (4.9,-1.7) {\begin{tabular}{c|cc}$C$&$P(R{=}T)$&$P(R{=}F)$\\\hline T&0.8&0.2\\F&0.2&0.8\end{tabular}};
\node[cpt,anchor=north east] at (-4.9,-3.0) {\begin{tabular}{c|cc}$C$&$P(S{=}T)$&$P(S{=}F)$\\\hline T&0.1&0.9\\F&0.5&0.5\end{tabular}};
\node[cpt,anchor=west] at (2.0,-4.8) {\begin{tabular}{cc|cc}$S$&$R$&$P(W{=}T)$&$P(W{=}F)$\\\hline T&T&0.99&0.01\\T&F&0.90&0.10\\F&T&0.90&0.10\\F&F&0.00&1.00\end{tabular}};
\end{tikzpicture}}
\caption{Mạng Bayes và các bảng CPT.}
\end{figure}

Phân rã xác suất đồng thời theo cấu trúc mạng:
\[
P(C,S,R,W)=P(C)\,P(S\mid C)\,P(R\mid C)\,P(W\mid S,R).
\]

\subsection*{a) $P(W=F\mid C=T)$}
Biên duyên hóa theo $S$ và $R$ (hai biến ẩn); do $W$ chỉ phụ thuộc $S,R$ và, khi biết $C$, $S$ độc lập với $R$:
\begin{align*}
P(W{=}F\mid C{=}T)&=\sum_{s\in\{T,F\}}\sum_{r\in\{T,F\}}P(W{=}F,S{=}s,R{=}r\mid C{=}T)\\
&=\sum_{s}\sum_{r}P(S{=}s\mid C{=}T)\,P(R{=}r\mid C{=}T)\,P(W{=}F\mid S{=}s,R{=}r).
\end{align*}
\begin{center}
\begin{tabular}{cc|c|c|c|c}
\toprule
$S$ & $R$ & $P(S\mid C{=}T)$ & $P(R\mid C{=}T)$ & $P(W{=}F\mid S,R)$ & Tích \\
\midrule
T & T & $0.1$ & $0.8$ & $0.01$ & $0.1\cdot0.8\cdot0.01=0.0008$ \\
T & F & $0.1$ & $0.2$ & $0.10$ & $0.1\cdot0.2\cdot0.10=0.0020$ \\
F & T & $0.9$ & $0.8$ & $0.10$ & $0.9\cdot0.8\cdot0.10=0.0720$ \\
F & F & $0.9$ & $0.2$ & $1.00$ & $0.9\cdot0.2\cdot1.00=0.1800$ \\
\midrule
\multicolumn{5}{r|}{Tổng} & $\mathbf{0.2548}$ \\
\bottomrule
\end{tabular}
\end{center}
\[
\boxed{P(W=F\mid C=T)=0.0008+0.0020+0.0720+0.1800=0.2548\ (25.48\%)}
\]
(Kiểm tra: $P(W{=}T\mid C{=}T)=0.08\cdot0.99+0.02\cdot0.9+0.72\cdot0.9+0.18\cdot0=0.7452$, và $0.7452+0.2548=1$.)

\subsection*{b) $P(R=T,S=T)$}
Biên duyên hóa theo $C$ (còn $W$ được lấy tổng bằng $1$ vì $\sum_w P(w\mid S,R)=1$):
\begin{align*}
P(R{=}T,S{=}T)&=\sum_{c\in\{T,F\}}\sum_{w\in\{T,F\}}P(c)\,P(S{=}T\mid c)\,P(R{=}T\mid c)\,P(w\mid S{=}T,R{=}T)\\
&=\sum_{c}P(c)\,P(S{=}T\mid c)\,P(R{=}T\mid c)\\
&=\underbrace{0.5\cdot0.1\cdot0.8}_{c=T}+\underbrace{0.5\cdot0.5\cdot0.2}_{c=F}=0.04+0.05.
\end{align*}
\[
\boxed{P(R=T,S=T)=0.09}
\]
\emph{Nhận xét:} $P(S{=}T)=0.5\cdot0.1+0.5\cdot0.5=0.3$ và $P(R{=}T)=0.5\cdot0.8+0.5\cdot0.2=0.5$, tích $=0.15\ne0.09$; vậy $S$ và $R$ \emph{không} độc lập khi không biết $C$ (do cùng phụ thuộc vào Cloudy).

%====================================================================
\clearpage
\cau{4}
\textbf{Đề:} Blocks World với $3$ khối A, B, C. Trạng thái đầu: B nằm trên A, A và C nằm trên bàn. Trạng thái đích: A nằm trên C, C và B nằm trên bàn.

\subsection*{a) Biểu diễn trạng thái}
Các vị từ tĩnh $\Blk{A},\Blk{B},\Blk{C}$ luôn đúng trong mọi trạng thái.
\begin{align*}
S_0\ (\text{đầu})=\{&\On{B}{A}\cm \OnT{A}\cm \OnT{C}\cm \Clr{B}\cm \Clr{C}\cm \HE\}\\
\mathcal{G}\ (\text{đích})=\{&\On{A}{C}\cm \OnT{C}\cm \OnT{B}\cm \Clr{A}\cm \Clr{B}\cm \HE\}
\end{align*}
(Ở cả hai trạng thái tay robot rỗng, tức $\HE$ đúng và không có $\Hold{X}$ nào đúng.)

\subsection*{b) Tiền điều kiện và tác động của các thao tác}
\begin{center}
\small
\renewcommand{\arraystretch}{1.35}
\begin{tabular}{@{}l >{\raggedright\arraybackslash}p{4.0cm} >{\raggedright\arraybackslash}p{3.7cm} >{\raggedright\arraybackslash}p{4.0cm}@{}}
\toprule
\textbf{Thao tác} & \textbf{Preconditions} & \textbf{Add-list} & \textbf{Delete-list} \\
\midrule
$\Op{UNSTACK}{X,Y}$ & $\On{X}{Y}\cm \Clr{X}\cm \HE$ & $\Hold{X}\cm \Clr{Y}$ & $\On{X}{Y}\cm \Clr{X}\cm \HE$ \\
$\Op{STACK}{X,Y}$ & $\Hold{X}\cm \Clr{Y}$ & $\On{X}{Y}\cm \Clr{X}\cm \HE$ & $\Hold{X}\cm \Clr{Y}$ \\
$\Op{PICKUP}{X}$ & $\OnT{X}\cm \Clr{X}\cm \HE$ & $\Hold{X}$ & $\OnT{X}\cm \Clr{X}\cm \HE$ \\
$\Op{PUTDOWN}{X}$ & $\Hold{X}$ & $\OnT{X}\cm \Clr{X}\cm \HE$ & $\Hold{X}$ \\
\bottomrule
\end{tabular}
\end{center}
(Trong mọi thao tác, $X,Y$ là các khối, $X\neq Y$.)

\subsection*{c) Lập kế hoạch bằng Goal Stack}
\paragraph{Nguyên tắc.} Khởi tạo ngăn xếp chỉ chứa mục tiêu $\mathcal{G}$, rồi lặp: xét phần tử ở đỉnh:
\begin{itemize}[nosep]
  \item \emph{Mục tiêu ghép} chưa thỏa: đẩy từng mệnh đề thành phần lên ngăn xếp (mục tiêu ghép ở dưới để kiểm tra lại sau cùng).
  \item \emph{Mệnh đề đơn} đã đúng trong trạng thái hiện tại: pop. Chưa đúng: thay bằng thao tác $\mathrm{op}$ có Add-list chứa mệnh đề đó, rồi đẩy $\mathrm{op}$, mục tiêu ghép các Preconditions của $\mathrm{op}$ và từng Precondition lên.
  \item \emph{Thao tác} ở đỉnh: pop, áp dụng vào trạng thái (xóa Delete-list, thêm Add-list) và ghi vào kế hoạch.
\end{itemize}
Các ký hiệu viết tắt:
$\mathcal{G}=\On{A}{C}\wedge\OnT{B}\wedge\OnT{C}\wedge\Clr{A}\wedge\Clr{B}\wedge\HE$;
$P_1=\PrA$ (precond. của STACK(A,C));
$P_2=\PrB$ (precond. của PICKUP(A));
$P_3=\PrC$ (precond. của UNSTACK(B,A)).
Ở mỗi bước, ngăn xếp được liệt kê từ \textbf{đỉnh} (trái) đến \textbf{đáy} (phải), sau khi thực hiện thao tác của bước đó.

\footnotesize
\begin{longtable}{@{}>{\centering\arraybackslash}p{0.6cm} >{\raggedright\arraybackslash}p{3.6cm} >{\raggedright\arraybackslash}p{9.7cm} >{\raggedright\arraybackslash}p{1.2cm}@{}}
\toprule
\textbf{Bước} & \textbf{Thao tác trên ngăn xếp} & \textbf{Nội dung ngăn xếp (đỉnh $\to$ đáy)} & \textbf{TT} \\
\midrule
\endhead
0 & Khởi tạo: push $\mathcal{G}$ & $\mathcal{G}$ & $S_0$ \\
1 & $\mathcal{G}$ chưa thỏa $\Rightarrow$ push các mệnh đề thành phần & $\On{A}{C}\sm \Rst$ & $S_0$ \\
2 & $\On{A}{C}$ sai $\Rightarrow$ pop; push $\Op{STACK}{A,C}$, $P_1$, $\Clr{C}$, $\Hold{A}$ & $\Hold{A}\sm \Clr{C}\sm P_1\sm \Op{STACK}{A,C}\sm \Rst$ & $S_0$ \\
3 & $\Hold{A}$ sai $\Rightarrow$ pop; push $\Op{PICKUP}{A}$, $P_2$, $\HE$, $\Clr{A}$, $\OnT{A}$ & $\OnT{A}\sm \Clr{A}\sm \HE\sm P_2\sm \Op{PICKUP}{A}\sm \Clr{C}\sm P_1\sm \Op{STACK}{A,C}\sm \Rst$ & $S_0$ \\
4 & $\OnT{A}$ đúng trong $S_0$ $\Rightarrow$ pop & $\Clr{A}\sm \HE\sm P_2\sm \Op{PICKUP}{A}\sm \Clr{C}\sm P_1\sm \Op{STACK}{A,C}\sm \Rst$ & $S_0$ \\
5 & $\Clr{A}$ sai $\Rightarrow$ pop; push $\Op{UNSTACK}{B,A}$, $P_3$, $\HE$, $\Clr{B}$, $\On{B}{A}$ & $\On{B}{A}\sm \Clr{B}\sm \HE\sm P_3\sm \Op{UNSTACK}{B,A}\sm \HE\sm P_2\sm \Op{PICKUP}{A}\sm \Clr{C}\sm P_1\sm \Op{STACK}{A,C}\sm \Rst$ & $S_0$ \\
6 & $\On{B}{A}\cm \Clr{B}\cm \HE\cm P_3$ đều đúng trong $S_0$ $\Rightarrow$ pop 4 phần tử & $\Op{UNSTACK}{B,A}\sm \HE\sm P_2\sm \Op{PICKUP}{A}\sm \Clr{C}\sm P_1\sm \Op{STACK}{A,C}\sm \Rst$ & $S_0$ \\
7 & pop và \textbf{áp dụng} $\Op{UNSTACK}{B,A}$ & $\HE\sm P_2\sm \Op{PICKUP}{A}\sm \Clr{C}\sm P_1\sm \Op{STACK}{A,C}\sm \Rst$ & $S_1$ \\
8 & $\HE$ sai (đang cầm B) $\Rightarrow$ pop; push $\Op{PUTDOWN}{B}$, $\Hold{B}$ & $\Hold{B}\sm \Op{PUTDOWN}{B}\sm P_2\sm \Op{PICKUP}{A}\sm \Clr{C}\sm P_1\sm \Op{STACK}{A,C}\sm \Rst$ & $S_1$ \\
9 & $\Hold{B}$ đúng trong $S_1$ $\Rightarrow$ pop & $\Op{PUTDOWN}{B}\sm P_2\sm \Op{PICKUP}{A}\sm \Clr{C}\sm P_1\sm \Op{STACK}{A,C}\sm \Rst$ & $S_1$ \\
10 & pop và \textbf{áp dụng} $\Op{PUTDOWN}{B}$ & $P_2\sm \Op{PICKUP}{A}\sm \Clr{C}\sm P_1\sm \Op{STACK}{A,C}\sm \Rst$ & $S_2$ \\
11 & $P_2=\OnT{A}\wedge\Clr{A}\wedge\HE$ đúng trong $S_2$ $\Rightarrow$ pop & $\Op{PICKUP}{A}\sm \Clr{C}\sm P_1\sm \Op{STACK}{A,C}\sm \Rst$ & $S_2$ \\
12 & pop và \textbf{áp dụng} $\Op{PICKUP}{A}$ & $\Clr{C}\sm P_1\sm \Op{STACK}{A,C}\sm \Rst$ & $S_3$ \\
13 & $\Clr{C}$ đúng, $P_1=\Hold{A}\wedge\Clr{C}$ đúng trong $S_3$ $\Rightarrow$ pop 2 phần tử & $\Op{STACK}{A,C}\sm \Rst$ & $S_3$ \\
14 & pop và \textbf{áp dụng} $\Op{STACK}{A,C}$ & $\Rst$ & $S_4$ \\
15 & $\OnT{B}\cm \OnT{C}\cm \Clr{A}\cm \Clr{B}\cm \HE$ đều đúng trong $S_4$; $\mathcal{G}$ thỏa $\Rightarrow$ pop hết & (rỗng) $\Rightarrow$ \textbf{HOÀN THÀNH} & $S_4$ \\
\bottomrule
\end{longtable}
\normalsize

\paragraph{Kế hoạch tìm được} ($4$ hành động):
\[
\Op{UNSTACK}{B,A}\ \to\ \Op{PUTDOWN}{B}\ \to\ \Op{PICKUP}{A}\ \to\ \Op{STACK}{A,C}.
\]

\paragraph{Trạng thái thế giới sau mỗi hành động} (bỏ qua các vị từ tĩnh $\Blk{\cdot}$):
\begin{center}
\small
\renewcommand{\arraystretch}{1.3}
\begin{tabular}{@{}l l >{\raggedright\arraybackslash}p{8.8cm}@{}}
\toprule
\textbf{Trạng thái} & \textbf{Hành động vừa áp dụng} & \textbf{Các mệnh đề đúng} \\
\midrule
$S_0$ & (ban đầu) & $\On{B}{A}\cm \OnT{A}\cm \OnT{C}\cm \Clr{B}\cm \Clr{C}\cm \HE$ \\
$S_1$ & $\Op{UNSTACK}{B,A}$ & $\OnT{A}\cm \OnT{C}\cm \Clr{A}\cm \Clr{C}\cm \Hold{B}$ \\
$S_2$ & $\Op{PUTDOWN}{B}$ & $\OnT{A}\cm \OnT{B}\cm \OnT{C}\cm \Clr{A}\cm \Clr{B}\cm \Clr{C}\cm \HE$ \\
$S_3$ & $\Op{PICKUP}{A}$ & $\OnT{B}\cm \OnT{C}\cm \Clr{B}\cm \Clr{C}\cm \Hold{A}$ \\
$S_4$ & $\Op{STACK}{A,C}$ & $\On{A}{C}\cm \OnT{B}\cm \OnT{C}\cm \Clr{A}\cm \Clr{B}\cm \HE$ \\
\bottomrule
\end{tabular}
\end{center}
$S_4$ chứa đầy đủ các mệnh đề của $\mathcal{G}$ nên kế hoạch đúng.

%====================================================================
\clearpage
\cau{5}
\textbf{Đề:} Phân loại Naive Bayes cho khái niệm \emph{Transportation Means}.

\begin{center}
\begin{tabular}{@{}cllllc@{}}
\toprule
No. & Gender & Cars Own & Transp. Cost & Income & Class \\
\midrule
1 & Male & 0 & Cheap & Low & Bicycle \\
2 & Male & 1 & Cheap & Average & Bicycle \\
3 & Female & 1 & Expensive & High & Car \\
4 & Male & 1 & Expensive & Average & Car \\
5 & Female & 2 & Expensive & High & Car \\
6 & Male & 0 & Normal & Low & Motorbike \\
7 & Female & 1 & Normal & Average & Motorbike \\
\bottomrule
\end{tabular}
\end{center}
Có $N=7$ mẫu: Bicycle $=2$ (mẫu 1, 2), Car $=3$ (mẫu 3, 4, 5), Motorbike $=2$ (mẫu 6, 7).

\subsection*{a) Prior và Likelihood}
\paragraph{Prior.}
\[
P(\text{Bicycle})=\frac27,\qquad P(\text{Car})=\frac37,\qquad P(\text{Motorbike})=\frac27 .
\]

\paragraph{Likelihood chưa làm mượt.} $P(x_i=v\mid C)=\dfrac{n_{v,C}}{n_C}$ với $n_C=2,3,2$ lần lượt cho Bicycle, Car, Motorbike. Nhiều giá trị bằng $0$ (Gender=Female$\mid$Bicycle, Cars=2$\mid$Bicycle, Cost=Normal$\mid$Bicycle, Income=High$\mid$Bicycle, \dots) nên cả \textbf{bốn thuộc tính} đều chứa xác suất $0$.

\paragraph{Laplace smoothing (add-1).} Với mỗi thuộc tính có ít nhất một xác suất bằng $0$, áp dụng cho \emph{mọi} giá trị của thuộc tính đó (ở mọi lớp):
\[
P(x_i=v\mid C)=\frac{n_{v,C}+1}{n_C+k_i},
\]
với $k_i$ là số giá trị khác nhau của thuộc tính ($k_{\text{Gender}}=2$, $k_{\text{Cars}}=3$, $k_{\text{Cost}}=3$, $k_{\text{Income}}=3$). Các mẫu số là $n_C+k_i$: Gender $\to 4,5,4$; Cars, Cost, Income $\to 5,6,5$ (Bicycle, Car, Motorbike).

\begin{center}
\small
\renewcommand{\arraystretch}{1.5}
\begin{tabular}{@{}ll ccc ccc@{}}
\toprule
 & & \multicolumn{3}{c}{Trước làm mượt $\dfrac{n}{n_C}$} & \multicolumn{3}{c}{Sau add-1 $\dfrac{n+1}{n_C+k}$} \\
\cmidrule(lr){3-5}\cmidrule(lr){6-8}
Thuộc tính & Giá trị & Bicycle & Car & Motorbike & Bicycle & Car & Motorbike \\
\midrule
Gender & Male   & $\frac22$ & $\frac13$ & $\frac12$ & $\frac34$ & $\frac25$ & $\frac24=\frac12$ \\
       & Female & $\frac02$ & $\frac23$ & $\frac12$ & $\frac14$ & $\frac35$ & $\frac24=\frac12$ \\
\midrule
Cars Own & 0 & $\frac12$ & $\frac03$ & $\frac12$ & $\frac25$ & $\frac16$ & $\frac25$ \\
         & 1 & $\frac12$ & $\frac23$ & $\frac12$ & $\frac25$ & $\frac36=\frac12$ & $\frac25$ \\
         & 2 & $\frac02$ & $\frac13$ & $\frac02$ & $\frac15$ & $\frac26=\frac13$ & $\frac15$ \\
\midrule
Transp. Cost & Cheap     & $\frac22$ & $\frac03$ & $\frac02$ & $\frac35$ & $\frac16$ & $\frac15$ \\
             & Expensive & $\frac02$ & $\frac33$ & $\frac02$ & $\frac15$ & $\frac46=\frac23$ & $\frac15$ \\
             & Normal    & $\frac02$ & $\frac03$ & $\frac22$ & $\frac15$ & $\frac16$ & $\frac35$ \\
\midrule
Income & Low     & $\frac12$ & $\frac03$ & $\frac12$ & $\frac25$ & $\frac16$ & $\frac25$ \\
       & Average & $\frac12$ & $\frac13$ & $\frac12$ & $\frac25$ & $\frac26=\frac13$ & $\frac25$ \\
       & High    & $\frac02$ & $\frac23$ & $\frac02$ & $\frac15$ & $\frac36=\frac12$ & $\frac15$ \\
\bottomrule
\end{tabular}
\end{center}
(Kiểm tra: mỗi cột sau làm mượt của một thuộc tính đều có tổng bằng $1$, ví dụ Cars$\mid$Car: $\tfrac16+\tfrac36+\tfrac26=1$.)

\subsection*{b) Phân loại hai mẫu mới}
Dùng $P(C\mid x)\propto P(C)\prod_i P(x_i\mid C)$ (giả thiết độc lập có điều kiện), với các likelihood đã làm mượt.

\paragraph{Mẫu 8: (Female, 0, Expensive, High).}
\begin{align*}
\text{Bicycle:}\quad & \tfrac27\cdot\tfrac14\cdot\tfrac25\cdot\tfrac15\cdot\tfrac15=\tfrac{4}{3500}=\tfrac{1}{875}\approx0.001143\\
\text{Car:}\quad & \tfrac37\cdot\tfrac35\cdot\tfrac16\cdot\tfrac23\cdot\tfrac12=\tfrac{1}{70}\approx0.014286\\
\text{Motorbike:}\quad & \tfrac27\cdot\tfrac12\cdot\tfrac25\cdot\tfrac15\cdot\tfrac15=\tfrac{2}{875}\approx0.002286
\end{align*}
Tổng (mẫu số chuẩn hóa) $=\tfrac{2}{1750}+\tfrac{25}{1750}+\tfrac{4}{1750}=\tfrac{31}{1750}\approx0.017714$. Xác suất hậu nghiệm:
\[
P(\text{Bicycle}\mid x_8)=\tfrac{2}{31}\approx0.0645,\quad
P(\text{Car}\mid x_8)=\tfrac{25}{31}\approx0.8065,\quad
P(\text{Motorbike}\mid x_8)=\tfrac{4}{31}\approx0.1290.
\]
\textbf{Kết luận:} mẫu 8 thuộc lớp \textbf{Car}.

\paragraph{Mẫu 9: (Male, 2, Normal, Low).}
\begin{align*}
\text{Bicycle:}\quad & \tfrac27\cdot\tfrac34\cdot\tfrac15\cdot\tfrac15\cdot\tfrac25=\tfrac{12}{3500}=\tfrac{3}{875}\approx0.003429\\
\text{Car:}\quad & \tfrac37\cdot\tfrac25\cdot\tfrac13\cdot\tfrac16\cdot\tfrac16=\tfrac{1}{630}\approx0.001587\\
\text{Motorbike:}\quad & \tfrac27\cdot\tfrac12\cdot\tfrac15\cdot\tfrac35\cdot\tfrac25=\tfrac{6}{875}\approx0.006857
\end{align*}
Tổng $=\tfrac{54}{15750}+\tfrac{25}{15750}+\tfrac{108}{15750}=\tfrac{187}{15750}\approx0.011873$. Xác suất hậu nghiệm:
\[
P(\text{Bicycle}\mid x_9)=\tfrac{54}{187}\approx0.2888,\quad
P(\text{Car}\mid x_9)=\tfrac{25}{187}\approx0.1337,\quad
P(\text{Motorbike}\mid x_9)=\tfrac{108}{187}\approx0.5775.
\]
\textbf{Kết luận:} mẫu 9 thuộc lớp \textbf{Motorbike}.

\begin{center}
\begin{tabular}{@{}lcccc@{}}
\toprule
Mẫu & $P(\text{Bicycle}\mid x)$ & $P(\text{Car}\mid x)$ & $P(\text{Motorbike}\mid x)$ & Dự đoán \\
\midrule
8: (Female, 0, Expensive, High) & $0.0645$ & $\mathbf{0.8065}$ & $0.1290$ & \textbf{Car} \\
9: (Male, 2, Normal, Low) & $0.2888$ & $0.1337$ & $\mathbf{0.5775}$ & \textbf{Motorbike} \\
\bottomrule
\end{tabular}
\end{center}

\end{document}
